\scanpage{222}
% opening system of data, lines 1–4
1) Soit $f:X\longrightarrow S$ morphisme lisse d'espaces analytiques (lisses)

$\omega$ une 1-forme fermée \inserted{(holomorphe)} relative sur $X/S$

$a$ et $b$ deux sections \inserted{holomorphes} de $X$ sur $S$, telles que $\forall s\in S$, $a(s)$ et $b(s)$ \inserted{soient dans la $\overline{\mathrm m}$ comp. connexe de $X_s$}

% local paths, lines 5–7
Localement sur $S$, on peut trouver un arc \deleted{\unsure{diff.|diff}} \inserted{(différentiable par morceaux)} $C_s:[0,1]=I\longrightarrow X_s$, joignant $a(s)$ et $b(s)$, et dépendant \unsure{continument|continûment} de $s$.

% primitive and the first integral, lines 8–12
Utilisant le fait que $\omega$ est fermée, donc \unsure{localement|localem\textsuperscript{t}} de la forme $d_{X/S}\varphi$, on constate que la fonction sur $S$
\[
 I(s)=\int_{\deleted{\unsure{a(s)}}C_s}^{\deleted{\unsure{b(s)}}}\omega_s
       =\int_{a(s)}^{b(s)}\omega_s
\]
est holomorphe en $s$. [\unsure{En|On|En effet} \inserted{\unsure{se plaçant|plus précisément} au voisinage de $s_0\in S$ et \unsure{se|en}} \unsure{ramenant|ramène|supposons} au cas où \deleted{\unsure{$a(s_0)$ et|$a(s_0)$}} $a(s_0)$ et $b(s_0)$ sont assez voisins pour que dans un ouvert $U$ contenant $a(s_0),b(s_0)$, on ait $\omega|U=d\varphi$, donc pour $s$ voisin de $s_0$, $I(s)=\varphi b(s)-\varphi a(s)$].

% change of paths and equation (1)
Changeant le système des $C_s$ revient à les remplacer par $C'_s=C_s+Z_s$, $Z_s$ un 1-cycle, donc on trouve
\[
 I'(s)=I(s)+\underbrace{\int_{Z_s}\omega_s}_{\deleted{\unsure{Z}}}
 \tag*{(1)}
\]

% bottom and equation (2)
Considérons une expression formelle (avec des $\omega_i$ comme $\omega$, et des op. diff. \inserted{(hol.)} $D_i$ sur $S$ ]
\[
 \sum D_{i\,s}\int_{a(s)}^{b(s)}\omega_{i\,s}
 \tag*{(2)}
\]
% archive sheet number: 221
